\newcommand{\tipoRequerimento}{cobrança indevida sobre enquadramento tarifário incorreto}
\section{DOS FATOS E DOS DIREITOS}

\begin{enumerate}[resume]
    \item O laudo de vistoria da distribuidora é nulo porque apenas alegou genericamente que a carga instalada não é exclusiva de saneamento, sem descrever ou identificar quais seriam as outras supostas cargas na unidade.
    
    \item A simples existência de tomadas não afasta o enquadramento na classe de Serviço Público. Essas tomadas são instaladas no interior de painéis para possibilitar a manutenção, calibração e operação do próprio sistema de saneamento, atuando como cargas de apoio essenciais.
    
    \item A distribuidora fundamentou o indeferimento alegando que havia um aparelho de choque ligado em uma das tomadas. No entanto, a Ordem de Serviço (OS) emitida no momento da vistoria não cita e não mostra esse aparelho, registrando apenas a presença de uma ``tomada/extensão''.
    
    \item Não há nenhuma foto ou registro técnico no processo que comprove a presença ou conexão de tal aparelho de choque. A alegação da concessionária é unilateral, tardia e sem qualquer prova material.
    
    \item Recusar o enquadramento tarifário com base em um equipamento que não foi listado ou fotografado pela equipe de campo no ato da vistoria invalida a fundamentação do indeferimento, violando a Resolução Normativa nº 1.000/2021 da ANEEL.
    
    \item Como a energia da unidade atende exclusivamente ao sistema de saneamento básico (bomba de poço artesiano), o enquadramento correto é a tarifa de Serviço Público.
    
    \item Diante da cobrança incorreta e da recusa injustificada em reclassificar a tarifa, o consumidor tem o direito de receber de volta os valores pagos a mais.
\end{enumerate}

\section{DOS PEDIDOS}

Diante do exposto, requer-se a esta ilustre Agência Reguladora:

\begin{enumerate}[resume]
    \item A intimação da concessionária para apresentar provas reais (como fotos com data) que comprovem a conexão do suposto aparelho de choque, já que a Ordem de Serviço original apontou apenas ``tomada/extensão'';
    
    \item A anulação da decisão de indeferimento, uma vez que foi baseada em alegações sem provas;
    
    \item A imediata correção do enquadramento tarifário da unidade consumidora nº 13896006 para a classe de Serviço Público, conforme a Resolução Normativa nº 1.000/2021 da ANEEL;
    
    \item A devolução dos valores cobrados a maior nos últimos 10 anos, conforme o Despacho nº 2.006 da ANEEL e o art. 205 do Código Civil;
    
    \item Que a devolução ocorra em dobro, corrigida monetariamente pelo IPCA e acrescida de juros de 1\% ao mês;
    
    \item Que o pagamento da devolução seja feito via depósito bancário na conta do Município, nos termos da REN nº 1.000/2021.
\end{enumerate}

